\section{A small-width Gaussian calculation for the coarse report}
\label{app:small-width}

The coarse report also admits a different Gaussian parameter choice.
We balance the interval width against the spacing of the selector grid
and use the same cancellation as in \eqref{eq:coarse-gaussian-gain}
to compute the resulting predictor error. We also show how this error
transfers when the two predictor constants are recalibrated to the
input law, using convergence of restricted moments.

\subsection{The parameters and a law-dependent error bound}
\label{app:small-width-rule}

For $0<w<1$, put
\begin{equation}\label{app:small-width-schedule}
 h=w^{2/3},\qquad T=\sqrt{4\log(1/w)}.
\end{equation}
List the points of $h\mathbb Z\cap[-T,T]$ in increasing order as
$t_1,\ldots,t_m$, and let $i_*$ index zero.  Thus
$m=2\lfloor T/h\rfloor+1=O(T/h)$ as $w\downarrow0$; in particular,
$m\ge3$ for sufficiently small $w$.  On $|t|\le T$, let $I(t)$
be the index of a nearest grid point, breaking ties by the smaller
index.  On $|t|>T$, set $I(t)=i_*$.  The grid endpoints need not
equal $\pm T$, but
\[
 |t-t_{I(t)}|\le h\qquad (|t|\le T).
\]
Give leaf $j$ the closed interval
\[
 A_j=[t_j+1-w/2,t_j+1+w/2].
\]
These endpoint and tie conventions define the rule on every real input.
Restrict henceforth to sufficiently small $w$ that $m\ge3$.

The choice $h=w^{2/3}$ balances two power scales: $h^2/w$ from
rounding the central input to a grid point, and $w/h$ from adding
the leaf-interval variances. The accompanying cutoff factors grow
more slowly than every power of $1/w$. The choice of $T$ makes
$e^{-T^2/2}=w^2$, so the Gaussian tail second moment is $o(w)$.
The estimates below retain all these factors.

Apply the coarse report of Lemma~\ref{lem:coarse-report} with central
input $t$, leaf inputs $u_1,\ldots,u_m$, selector $I$, and leaf
intervals $A_j$.  Denote it by $\mathcal O_w$.  Thus it returns $Y$
when all leaves pass, $X$ when only the selected leaf fails, and
$(O,t,I(t),(u_j)_{j\ne I(t)})$ otherwise.  In the notation of that
lemma, these events are $B$, $A=P\setminus B$, and $P^c$.
For clarity, the special-cell identity here is
\begin{equation}\label{app:small-width-cell}
 \one_A(t,u)=
 \sum_{i=1}^m\one_{\{I(t)=i\}}\prod_{j\ne i}\one_{A_j}(u_j)
       -\prod_{j=1}^m\one_{A_j}(u_j).
\end{equation}
The lemma proves that every cell has coordinate cost at most $m$:
the displayed summands use $m$ coordinates, the $Y$ cell ignores $t$,
and an ordinary report imposes no condition on its omitted leaf.
This holds pointwise on every finite product alphabet, including
unused labels.  The boundaries of $A$, $B$, and $P$ lie in finitely
many coordinate hyperplanes.

Let $U$ be centered with variance one, and take independent copies
$U_0,\ldots,U_m$.  Write
\[
 L=\sum_{j=0}^mU_j,\qquad
 \mathcal R_w(U)=\E[(L-\E[L\mid\mathcal O_w(\mathbf U)])^2].
\]
For the moment assume
\[
 q_j=\Pp(U\in A_j)\in(0,1),\qquad Q=\prod_{j=1}^m q_j.
\]
Let $\mu_j,a_j$ be the conditional mean and variance on $A_j$,
and put
\[
 M_\mu=\sum_j\mu_j,\qquad S_a=\sum_j a_j.
\]
The interval length gives $0\le a_j\le w^2$.
Use the predictor $M_\mu-1$ on $X$, the predictor $M_\mu$ on $Y$,
and the sum of reported coordinates on ordinary outputs. Denote this
function of the report by $\widehat L$. We will compute its deficit
$1-\E[(L-\widehat L)^2]$ exactly.

The ordinary reports contribute $\Pp(P^c)$ to the squared error,
by Lemma~\ref{lem:coarse-report}. For the exceptional reports, fix
$U_0=t$ and put $i=I(t)$. The event $P$ has conditional probability
$Q/q_i$, restricts every unselected leaf to its interval, and leaves
$U_i$ unrestricted. Under these restrictions, $L-(M_\mu-1)$ has mean
$t+1-\mu_i$ and variance $1+S_a-a_i$. Hence
\[
 \E\bigl[\one_P(1-(L-(M_\mu-1))^2)\bigr]
 =-Q\E_t\left[\frac{(t+1-\mu_i)^2+S_a-a_i}{q_i}\right].
\]
On $B$, the central input remains unrestricted. The two respective
predictor deficits on $B$ are therefore
\begin{align*}
 \E\bigl[\one_B(1-(L-(M_\mu-1))^2)\bigr]&=-Q(1+S_a),\\
 \E\bigl[\one_B(1-(L-M_\mu)^2)\bigr]&=-QS_a.
\end{align*}
Subtracting the first of these from the $P$ contribution gives the
$X$ contribution, because $A=P\setminus B$. Adding the $Y$
contribution cancels $S_a$, exactly as in the earlier Gaussian
calculation. Define
\begin{equation}\label{app:small-width-deficit}
 H_w(U)=1-\E_t\left[
       \frac{(t+1-\mu_{I(t)})^2+S_a-a_{I(t)}}{q_{I(t)}}\right].
\end{equation}
We have proved
\begin{equation}\label{app:small-width-residual}
 1-\mathcal R_w(U)\ge
 1-\E[(L-\widehat L)^2]=QH_w(U).
\end{equation}
The equality is the gain of the specified predictor; the inequality
uses optimality of conditional expectation. The calculations integrate
against the law of $t$, so they require no positive probability for
any individual central value. We next show that both subtracted
errors in \eqref{app:small-width-deficit} tend to zero for Gaussian
inputs. The product $Q$ may be extremely small; only its strict
positivity will matter after $w$ is fixed.

\subsection{Gaussian estimates, including the rare-probability factors}
\label{app:small-width-estimates}

In this subsection $U=G$ is standard normal and $w\downarrow0$.
Write $\phi$ for its density.  Uniformly over the grid,
\[
 q_j\le\frac{w}{\sqrt{2\pi}},\qquad
 |\mu_j-(t_j+1)|\le w/2,\qquad
 \sum_ja_j\le mw^2=O(Tw^{4/3})\longrightarrow0.
\]
We must control the two errors in \eqref{app:small-width-deficit}
after division by the selected interval probability. A uniform lower
bound of order $w$ for all $q_j$ is unavailable, because the grid
extends into the Gaussian tails. Instead we compare each interval's
density with the density of the central coordinate that selects it.

For $|t|\le T$ and $z\in A_{I(t)}$, we have
$|z-t|\le1+h+w/2$.  Hence, for sufficiently small $w$,
\[
 q_{I(t)}\ge w\inf_{z\in A_{I(t)}}\phi(z),\qquad
 \frac{\phi(t)}{q_{I(t)}}
 \le w^{-1}\exp(C(1+T)).
\]
The constant $C$ is absolute: the logarithm of the density ratio is
$(z^2-t^2)/2\le |t||z-t|+|z-t|^2/2$.
Since $T=2\sqrt{\log(1/w)}$, for each fixed $\varepsilon>0$,
\begin{equation}\label{app:small-width-inverse-density}
 \frac{\phi(t)}{q_{I(t)}}=O(w^{-1-\varepsilon})
 \qquad (|t|\le T).
\end{equation}
The fallback interval is centered at $1$, so $q_{i_*}\ge cw$ for
an absolute $c>0$ and small $w$.  Integrating the central bound and
using this fallback bound on the tails gives
\begin{equation}\label{app:small-width-inverse-mean}
 \E\frac1{q_{I(G)}}=O(Tw^{-1-\varepsilon}).
\end{equation}
Take $\varepsilon=1/6$. Since $0\le S_a-a_{I(G)}\le mw^2$,
the variance contribution in
\eqref{app:small-width-deficit} is at most
\[
 mw^2\E\frac1{q_{I(G)}}
 =O(T^2w^{1/6})\longrightarrow0.
\]

For the squared mismatch, first suppose $|t|\le T$. Then
\[
 |t+1-\mu_{I(t)}|
 \le |t-t_{I(t)}|+w/2\le h+w/2.
\]
Its central contribution is therefore at most
\[
 O\bigl(Tw^{-7/6}(h+w)^2\bigr)
 =O(Tw^{1/6})\longrightarrow0.
\]
On $|t|>T$, the fallback mean is bounded, so the contribution is
\[
 O\left(w^{-1}\int_{|t|>T}(1+t^2)\phi(t)\,dt\right)
 =O\bigl(w^{-1}(1+T)e^{-T^2/2}\bigr)
 =O((1+T)w)\longrightarrow0.
\]
For the integral estimate, integration by parts gives
$\int_T^\infty t^2\phi(t)\,dt
=T\phi(T)+\int_T^\infty\phi(t)\,dt$, and
$\int_T^\infty\phi(t)\,dt\le\phi(T)/T$ for $T>0$.
The last equality uses $e^{-T^2/2}=w^2$.
Both subtracted errors therefore vanish, proving
\begin{equation}\label{app:small-width-limit}
 H_w(G)\longrightarrow1.
\end{equation}
Thus the same coarse report has a strictly positive Gaussian
improvement for this independent choice of parameters.  To use it
repeatedly on finite cubes, we now fix those parameters and transfer
the particular lower bound just proved.

\subsection{Restricted moments and the finite-score transfer}
\label{app:small-width-transfer}

Fix a sufficiently small $w>0$ with $m\ge3$ and $H_w(G)>1/2$.
All intervals, grid points, selector cells, and $m$ are fixed from
now on.  Put
\[
 Q_G=\prod_{j=1}^m\Pp(G\in A_j)>0,\qquad \eta=Q_G/4.
\]
Let $U_n$ be centered variance-one variables such that
$U_n\Rightarrow G$ and $\sup_n\E U_n^4<\infty$.
Write $U_{n,0},\ldots,U_{n,m}$ for independent copies.  Use a
subscript $n$ for the interval moments of $U_n$ and a superscript $G$
for their Gaussian values, and put $Q_n=\prod_j q_{j,n}$.
For each bounded interval $A_j$, weak convergence off its Gaussian-null
boundary gives convergence of its probability and its restricted
first and second moments.  Thus
\[
 q_{j,n}\to q_j^G\in(0,1),\qquad
 \mu_{j,n}\to\mu_j^G,\qquad a_{j,n}\to a_j^G.
\]
All these conditional moments are defined for sufficiently large $n$.
There are finitely many intervals, so their probabilities are then
simultaneously bounded away from zero and one.  In particular, every
coefficient in \eqref{app:small-width-deficit} converges, including
the reciprocal-probability factors.

Coefficient convergence alone does not yet transfer its expectation.
Partition the line into the fixed selector cells
$E_i=\{t:I(t)=i\}$, whose boundaries are finite.  The fallback cell
includes both tails as well as the central cell for the grid point
zero; it still has a finite boundary.  For $k=0,1,2$,
\begin{equation}\label{app:small-width-restricted-moments}
 \E[U_n^k\one_{\{U_n\in E_i\}}]
 \longrightarrow\E[G^k\one_{\{G\in E_i\}}].
\end{equation}
To see this, multiply $t^k\one_{E_i}(t)$ by a continuous cutoff that
is one on $[-K,K]$ and zero outside $[-K-1,K+1]$.  The resulting
function is bounded and has only Gaussian-null discontinuities, so
weak convergence applies.  The cutoff errors for $k=1,2$ vanish
uniformly as $K\to\infty$, by
\[
 \E[|U_n|\one_{\{|U_n|>K\}}]\le K^{-1},\qquad
 \E[U_n^2\one_{\{|U_n|>K\}}]
 \le K^{-2}\sup_j\E U_j^4.
\]
The corresponding Gaussian errors also vanish.  The case $k=0$
follows directly from the continuity-set criterion.  These are the
same weak-convergence and uniform-integrability principles used in
Lemma~\ref{lem:transfer}; see~\cite{Durrett}.

On a fixed selector cell $E_i$, the integrand in
\eqref{app:small-width-deficit} is a polynomial of degree at most two
in $t$. Expanding that polynomial expresses its expectation as a
finite sum of convergent coefficients times the restricted moments
in \eqref{app:small-width-restricted-moments}.  Therefore
\[
 Q_nH_w(U_n)\longrightarrow Q_GH_w(G)>Q_G/2.
\]
This proves convergence despite the moving predictor coefficients;
we have not assumed continuity of conditional variances.
For all sufficiently large $n$, \eqref{app:small-width-residual}
now gives $1-\mathcal R_w(U_n)>\eta$.  Since the sum of the
$m+1$ independent inputs is centered and has variance $m+1$,
\[
 \Var\!\left(\E\!\left[\sum_{j=0}^mU_{n,j}
                  \,\middle|\,\mathcal O_w(\mathbf U_n)\right]\right)
 =m+1-\mathcal R_w(U_n)>m+\eta.
\]
The pointwise coordinate-cost bound from
Lemma~\ref{lem:coarse-report}, with the substitution specified above,
and this transferred gain verify both hypotheses of
Proposition~\ref{prop:amplification} with
$(r,d,\eta)=(m+1,m,Q_G/4)$. The criterion gives the signed
conclusion of Theorem~\ref{thm:main}. Each averaging limit uses one
fixed finite observation and fixed parameters.  At that stage, the normalized
batch scores have a fourth-moment bound uniform in the batch size,
by \eqref{eq:fourth}; no uniformity across subsequent stages is required.
